\documentclass[aps,pre,reprint,amsmath,amssymb,nofootinbib,longbibliography,floatfix]{revtex4-2}
\usepackage[T1]{fontenc}
\usepackage{graphicx}
\usepackage{booktabs}
\usepackage{microtype}
\usepackage{tikz}
\usepackage[colorlinks=true,allcolors=blue!50!black]{hyperref}
\newcommand{\Tc}{T_{\mathrm c}}
\newcommand{\qm}{q_{\max}}
\newcommand{\atanh}{\operatorname{artanh}}
\newcommand{\Fan}{\Lambda_N}
\hypersetup{pdftitle={A periodic planar counterexample to the Apollonian Ising temperature bound},pdfauthor={Xiaojian Zeng}}
\begin{document}
\title{A periodic planar counterexample to the Apollonian Ising temperature bound}
\author{Xiaojian Zeng}
\email{xiaojianzeng@foxmail.com}
\affiliation{PBC School of Finance, Tsinghua University, Beijing 100190, China}
\date{October 14, 2026}
\begin{abstract}
An Apollonian continuation has recently been conjectured to bound the critical temperature of every periodic planar Ising ferromagnet with uniform exchange and given maximum coordination. We give an explicit counterexample within the stated Euclidean graph class. Retaining the triangular lattice, we insert a finite nested fan of 1024 spins into each upward triangle, with the same exchange \(J>0\) on every original and added bond and zero field. The resulting simple periodic triangulation has 1025 sites per primitive cell and maximum degree 2055. Exact local spin summation, together with the signed-homology characterization of the infinite-volume transition, proves \(\Tc/J>2/\ln(11/10)>20\). Monotonicity and exact interpolation of the source-defined canonical function give \(\Tc^*(2055)/J\leq\Tc^*(3072)/J<20\), using only nine finite radical steps. The separation disproves the universal Apollonian ceiling without changing the bare exchange, using non-Euclidean geometry, or taking an infinite decoration limit. The established full-face triangulation results and the weaker proved coordination bound retain their validity.
\end{abstract}
\maketitle
\raggedbottom

\section{Introduction}
The equilibrium ordering temperature of an Ising ferromagnet depends on the arrangement of its interaction graph. For uniform exchange \(J\), a bound involving the largest coordination number \(\qm\) would constrain how much resistance to thermal disorder can be obtained by rewiring a planar network. Joseph and Boettcher proved the bound
\begin{equation}
 \tanh(J/\Tc)\geq\tan\!\left(\frac{\pi}{2\qm}\right)
 \label{eq:knownbound}
\end{equation}
for the two-dimensional model in their stated class~\cite{JosephBoettcher2026}.

Joseph, Walsh, and Boettcher subsequently solved families generated by inserting a three-spoke spin into every triangular face at each iteration~\cite{JosephWalshBoettcher2026}. Their critical temperatures grow logarithmically with maximum coordination. The family starting from the triangular lattice defines an exact canonical continuation \(\Tc^*(q)\). They conjectured that
\begin{equation}
 \Tc(\Lambda)\leq\Tc^*(\qm(\Lambda))
 \label{eq:conjecture}
\end{equation}
for every periodic tessellation of the Euclidean plane in the uniform, zero-field model, with \(\qm\geq6\).

We disprove Eq.~\eqref{eq:conjecture} with one finite periodic triangulation. The construction concentrates distinct short ferromagnetic paths inside selected faces instead of subdividing every face at each step. The thermodynamic conclusion uses the periodic planar criticality theorem of Cimasoni and Duminil-Copin~\cite{CimasoniDuminilCopin2013}, applied to the original graph. Local decoration and periodic Ising criticality are established methods; the contribution is a strict counterexample to the proposed ceiling. We do not determine a replacement optimum or the smallest counterexample.

\section{Uniform model and finite planar witness}
Let
\begin{equation}
 \boldsymbol e_1=(1,0),\qquad
 \boldsymbol e_2=(1/2,\sqrt3/2),
\end{equation}
and place a hub spin at every \(\boldsymbol z=z_1\boldsymbol e_1+z_2\boldsymbol e_2\), with \(z_1,z_2\in\mathbb Z\). Retain all triangular-grid bonds. An upward triangle has corners \(A=\boldsymbol z\), \(B=\boldsymbol z+\boldsymbol e_1\), and \(C=\boldsymbol z+\boldsymbol e_2\). For fixed finite \(N\), insert
\begin{equation}
 \boldsymbol x_j=\boldsymbol z+
       (1/2,\sqrt3/(6j)),\qquad 1\leq j\leq N.
 \label{eq:coordinates}
\end{equation}
Add the bonds \(A x_j\) and \(B x_j\) for all \(j\), the bond \(C x_1\), and \(x_{j-1}x_j\) for \(2\leq j\leq N\). Repeat this rule in every upward triangle. Downward triangles are unchanged. Figure~\ref{fig:fan} illustrates the local construction.

\begin{figure*}[t]
 \centering
 \input{manuscript/figures/fan.tex}
 \caption{A single upward face with the nested fan at \(N=6\), drawn from Eq.~\eqref{eq:coordinates}. Black edges are retained triangular bonds; blue edges are added bonds. All have the same bare exchange \(J\). Each new spin lies strictly inside the remaining triangle \(ABx_{j-1}\); the actual witness uses fixed \(N=1024\). Only upward faces are decorated in the periodic lift. The original vector drawing uses AI-assisted figure code and represents the coordinate rule, rather than a simulated spin state.}
 \label{fig:fan}
\end{figure*}

The model is
\begin{equation}
 H=-J\sum_{\{i,j\}\in E(\Fan)}s_i s_j,\qquad
 s_i=\pm1,\qquad J>0,
 \label{eq:hamiltonian}
\end{equation}
with each distinct undirected edge counted once and no field. Set \(k_{\mathrm B}=1\),
\begin{equation}
 K=J/T,\quad r=e^{2K},\quad t=\tanh K=\frac{r-1}{r+1}.
 \label{eq:variables}
\end{equation}
Both \(J\) and \(T\) have energy units. The exchange is attached to graph adjacency, as in Ref.~\cite{JosephWalshBoettcher2026}, rather than normalized by degree or Euclidean bond length.

The first inserted site is the centroid of \(ABC\). For \(j\geq2\), its barycentric coordinates in \(A,B,x_{j-1}\) are
\begin{equation}
 \left(\frac1{2j},\frac1{2j},\frac{j-1}{j}\right).
\end{equation}
They are strictly positive. Each insertion is therefore an interior stellar subdivision of a triangular face. Induction proves noncrossing straight edges and positive-area triangular faces covering the original triangle. Repetition over disjoint face interiors gives a Euclidean triangulation of the plane.

For fixed finite \(N\), the cell contains finitely many distinct sites with positive minimum separation. The infinite graph is simple, locally finite, connected, and periodic under the two independent translations. Its cell counts are
\begin{equation}
 V=N+1,\quad E=3(N+1),\quad F=2(N+1).
 \label{eq:counts}
\end{equation}
A hub is an \(A\), a \(B\), and a \(C\) corner of one decorated face each. Its degree is \(6+N+N+1=2N+7\). For \(N\geq2\), the first \(N-1\) internal sites have degree four and the last degree three. Thus
\begin{equation}
 N=1024:\quad V=1025,\quad \qm=2055.
 \label{eq:witness}
\end{equation}
The average degree remains six. The triangular hub set and the different internal degrees identify the stated primitive translations. No limit \(N\to\infty\) is used to define the witness.

\section{Exact trace and thermodynamic ordering}
\subsection{Summing the physical fan}
For boundary spins \(A,B,C=\pm1\), sum only the \(N\) internal spins:
\begin{equation}
\begin{split}
 W_N(A,B,C)=\sum_{\boldsymbol x\in\{\pm1\}^N}
 \exp\!\bigg\{K\bigg[
  (A+B)\sum_{j=1}^N x_j\\
  +C x_1+\sum_{j=1}^{N-1}x_jx_{j+1}\bigg]\bigg\}.
\end{split}
\label{eq:trace}
\end{equation}
These are exactly the added uniform bonds; the three boundary bonds remain outside the trace. If \(A=-B\), summation from the free chain end gives \(W_0=(2\cosh K)^N\), independently of \(C\). Define
\begin{equation}
 U_+=\frac{W_N(+,+,+)}{W_0},\qquad
 U_-=\frac{W_N(+,+,-)}{W_0}.
\end{equation}
Reflection of \(A,B\) and global spin reversal then give, for every boundary configuration,
\begin{equation}
 W_N(A,B,C)=c\,e^{L_h AB+L_v(AC+BC)},
 \label{eq:localidentity}
\end{equation}
where
\begin{equation}
\begin{aligned}
 L_h&=\tfrac14\ln(U_+U_-),&
 L_v&=\tfrac14\ln(U_+/U_-),\\
 c&=W_0e^{L_h}>0.
\end{aligned}
\end{equation}
These effective parameters are a representation of the original uniform Hamiltonian.

Splitting the \(2K\) internal fields between neighboring factors in Eq.~\eqref{eq:trace} gives the rational transfer
\begin{equation}
 M=\frac1{r+1}
 \begin{pmatrix}r^2&1\\1&1\end{pmatrix},\qquad
 v=\frac1{r+1}\binom r1,
\end{equation}
and
\begin{equation}
 U_+=(r,r^{-1})M^{N-1}v,\qquad
 U_-=(1,1)M^{N-1}v.
 \label{eq:transfer}
\end{equation}
This follows by direct summation of Eq.~\eqref{eq:trace}, without an approximation in \(N\) or \(K\).

For \(r>1\), each vector \(w=M^k v\) has \(w_+>w_->0\). Its component ratio maps to \((r^2z+1)/(z+1)>1\). The sum of its components increases, since the increment is \((r-1)(rw_+-w_-)/(r+1)>0\). Consequently \(U_-\geq1\), and \(U_+-U_-=(r-1)(w_+-w_-/r)>0\). Both \(L_h,L_v\) are positive. Retaining the old edges produces one horizontal effective coupling \(K_h=K+L_h\) and two oblique couplings \(K_v=K+L_v\).

\subsection{The original graph's phase sign}
For a periodic planar ferromagnet, let \(C_h\) be the sum of weights of even edge subsets of a torus cell in mod-two homology class \(h\). The relevant signed quantity is
\begin{equation}
 F=C_{00}-C_{10}-C_{01}-C_{11}.
 \label{eq:signed}
\end{equation}
With positive original exchanges, the theorem in Ref.~\cite{CimasoniDuminilCopin2013} identifies its unique physical zero with the infinite-volume critical inverse temperature: \(F>0\) below that inverse temperature and \(F<0\) in the ordered phase. The graph above satisfies the theorem's local finiteness, periodicity, planarity, and disk-face hypotheses.

On the one-hub triangular quotient, the three old loops have translation labels \((1,0),(0,1),(-1,1)\). The empty and all-three subsets have trivial mod-two homology; the other six do not. With weights \(a=\tanh K_h\), \(b=\tanh K_v\), Eq.~\eqref{eq:signed} is
\begin{equation}
 P(a,b)=1-a-2b-2ab-b^2+ab^2.
 \label{eq:triangleP}
\end{equation}
Let \(X=e^{2K_h}=r\sqrt{U_+U_-}\) and \(Y=e^{2K_v}=r\sqrt{U_+/U_-}\). Direct substitution yields
\begin{equation}
 P(a,b)=
 \frac{4[2Y-X(Y^2-1)]}{(X+1)(Y+1)^2}.
\end{equation}
All denominators are positive, and therefore
\begin{equation}
 P(a,b)<0\quad\Longleftrightarrow\quad r^2U_+-U_->2.
 \label{eq:criterion}
\end{equation}

The local trace preserves the sign of the \emph{original} graph's quantity. Apply Eq.~\eqref{eq:localidentity} to each of the four flat torus boundary twists. Its restriction to a contractible decorated face is a vertex gauge, so the spin-independent factor \(c\) is unchanged and the same twist remains on the retained edges. The high-temperature expansion writes the four twisted partition functions as the four Fourier combinations of their even-subgraph homology sectors. Fourier inversion therefore multiplies every sector by the same positive normalization. For the primitive cell of \(\Fan\), explicit normalization gives
\begin{equation}
 F_N\!\left(\frac{r-1}{r+1}\right)
 =\frac{2^{2N+2}r^{N+1}}{(r+1)^{2N+3}}
       (2-r^2U_++U_-).
 \label{eq:originalF}
\end{equation}
This applies even when face boundary vertices are identified in the quotient. It certifies the phase of the original fixed-exchange model; no temperature-dependent effective Hamiltonian is substituted into the thermodynamic theorem.

\subsection{One finite exact certificate}
Take \(r_0=11/10\), hence \(K_0=\ln(11/10)/2\) and \(t_0=1/21\). Equation~\eqref{eq:transfer} can be evaluated with integers alone. Start at \((w_1,w_2)=(11,10)\) and apply
\begin{equation}
 (w_1,w_2)\longmapsto
 (121w_1+100w_2,\;100w_1+100w_2)
 \label{eq:integer}
\end{equation}
1023 times. With \(D=21\cdot210^{1023}\),
\begin{equation}
 U_+=\frac{121w_1+100w_2}{110D},\qquad
 U_-=\frac{w_1+w_2}{D}.
\end{equation}
The exact finite integer inequality is
\begin{equation}
 3641w_1+1100w_2>33000D.
 \label{eq:integercertificate}
\end{equation}
It implies \(r_0^2U_+-U_--2>1>0\). Appendix~\ref{app:arithmetic} supplies a short integer verification; the accompanying code retains the complete rational certificate. Its decimal margin is about \(34.36556\), for orientation only.

Equations~\eqref{eq:originalF} and \eqref{eq:integercertificate} give \(F_{1024}(1/21)<0\). The original lattice is strictly ordered at \(K_0\), so \(K_{\mathrm c}<K_0\). It follows that
\begin{equation}
 \boxed{\frac{\Tc(\Lambda_{1024})}{J}>
       \frac2{\ln(11/10)}>20.}
 \label{eq:lower}
\end{equation}
The second inequality follows from \(\ln(1+x)<x\) for \(x>0\). This is an infinite-volume phase bound, with no large-cell critical-root estimate or finite-size extrapolation.

\section{Comparison with the exact canonical function}
It is essential to compare with the function defined in Ref.~\cite{JosephWalshBoettcher2026}, rather than its rounded asymptotic constant. Write
\begin{equation}
\begin{aligned}
 g(t)&=\frac{2t}{1+t+\sqrt{1+6t-7t^2}},\\
 h(t)&=g^{-1}(t)=\frac{t(1+t)}{1-t+2t^2},\\
 A_0&=\frac2{\ln2},\qquad t_\triangle=2-\sqrt3.
\end{aligned}
\end{equation}
The source defines
\begin{equation}
 \kappa(t)=\lim_{n\to\infty}
       \left[\frac1{g^n(t)}-2(n-\ln n)\right],
\end{equation}
and
\begin{equation}
 C_\triangle=A_0\ln6-\kappa(t_\triangle)-2\ln(\ln2).
\end{equation}
For \(q\geq6\), put
\begin{equation}
 z_n(q)=A_0\ln(2^nq)-2\ln\!\ln(2^nq)-C_\triangle.
\end{equation}
The exact continuation is
\begin{equation}
 t^*(q)=\lim_{n\to\infty}h^n(1/z_n(q)),\qquad
 \frac{\Tc^*(q)}J=\frac1{\atanh t^*(q)}.
 \label{eq:canonicaldefinition}
\end{equation}
The source establishes the required limits. Iterates mean functional composition.

For sufficiently large \(n\), \(z_n(q)>1\) and
\begin{equation}
 z'_n(q)=\frac{A_0-2/\ln(2^nq)}q>0.
\end{equation}
Also \(h'(t)=(1-t)(1+3t)/(1-t+2t^2)^2>0\) on \((0,1)\). Taking the defining limits shows that \(t^*(q)\) is nonincreasing and \(\Tc^*(q)\) nondecreasing. In particular,
\begin{equation}
 \Tc^*(2055)\leq\Tc^*(3072).
 \label{eq:monotonecomparison}
\end{equation}

The triangular seed and doubling interpolation follow exactly from the definition. Let \(y_n=1/g^n(t_\triangle)\). By the definition of \(\kappa\),
\begin{equation}
 y_n=2(n-\ln n)+\kappa(t_\triangle)+o(1),
\end{equation}
whereas
\begin{equation}
\begin{split}
 z_n(6)=2(n-\ln n)+\kappa(t_\triangle)\\
 -2\ln\!\left(1+\frac{\ln6}{n\ln2}\right).
\end{split}
\end{equation}
Thus \(z_n(6)-y_n\to0\). The reciprocal forward map
\begin{equation}
 \eta(u)=\frac1{h(1/u)}=u-2+\frac4{u+1}
\end{equation}
maps \((1,\infty)\) into itself and has \(0<\eta'(u)<1\). Its iterates are nonexpanding, giving
\begin{equation}
 \left|\eta^n(z_n(6))-\eta^n(y_n)\right|
 \leq |z_n(6)-y_n|\longrightarrow0.
\end{equation}
Since \(\eta^n(y_n)=1/t_\triangle\), this proves \(t^*(6)=t_\triangle\).
Next, \(z_n(2q)=z_{n+1}(q)\). Reindexing the limits and using continuity of \(h\) gives
\begin{equation}
 t^*(2q)=g(t^*(q)),\qquad
 t^*(3072)=g^9(2-\sqrt3).
 \label{eq:interpolation}
\end{equation}

The comparison is a finite radical expression. Its reciprocal obeys
\begin{equation}
\begin{aligned}
 y_0&=2+\sqrt3,\\
 y_{j+1}&=\frac{y_j+1+\sqrt{y_j^2+6y_j-7}}2,
 \qquad 0\leq j<9.
\end{aligned}
\label{eq:radicals}
\end{equation}
Outward rational square-root bounds, with each squared endpoint checked exactly, give \(y_9<20\) (its display value is \(18.4410946764\ldots\)). The recurrence is increasing for \(y>1\), so endpoint propagation preserves the enclosure. Appendix~\ref{app:arithmetic} specifies the finite arithmetic; the exact endpoints are supplied with the code. Since \(\atanh t>t\) for \(0<t<1\),
\begin{equation}
 \boxed{\frac{\Tc^*(2055)}J\leq
        \frac{\Tc^*(3072)}J<y_9<20.}
 \label{eq:upper}
\end{equation}
The infinite defining limit is used analytically for monotonicity and interpolation. It is never numerically truncated.

\section{Checks, physical meaning, and scope}
For \(N=1\), the construction is the Half-Interpolated Triangular lattice of Ref.~\cite{JosephWalshBoettcher2026}, with \(\qm=9\). Its trace gives \(U_+=r+r^{-1}-1\) and \(U_-=1\). The transition is the unique physical root of
\begin{equation}
 r^3-r^2+r-3=0,
\end{equation}
whose derivative \(3r^2-2r+1\) is positive. Rational root isolation and a bounded logarithm series reproduce \(\Tc/J=4.4043925499\ldots\), agreeing with the reported \(4.404\). The undecorated triangular sign is \((1+t)(1-4t+t^2)\), giving \(t_{\mathrm c}=2-\sqrt3\) and \(\Tc/J=4/\ln3\). The full-face star trace recovers the source map \(h\) and its Laves-Star benchmark \(\Tc/J=5.00704969054\ldots\). These are calibrations against established results.

The finite ordering certificate was checked in three representations: the symmetric transfer in Eq.~\eqref{eq:transfer}, a recursion multiplying the original satisfied-bond factors, and direct elimination of original even-edge subsets while retaining both mod-two winding bits. The last calculation processes all 3075 witness edges, independently of the effective-coupling trace. At small \(N\), complete subset enumeration checks the homology elimination. A geometric weighted Kac--Ward calculation also gives \(\det(I-W)=P(a,b)^2\) symbolically. Two rational implementations enclose the canonical radicals, one in \(t\) and one in \(1/t\). All decisive operations are exact integer, rational, or analytic operations.

Combining Eqs.~\eqref{eq:lower} and \eqref{eq:upper} gives a strict violation of Eq.~\eqref{eq:conjecture} by an admissible finite-cell lattice. Many distinct two-edge paths enhance the effective horizontal interaction, while positive transverse interactions sustain bulk order. Those paths share endpoints but contain distinct spins and bonds. Every physical bond still has the same \(J\).

The witness has unequal geometric bond lengths and narrow triangular faces. Both are permitted by the source's graph model, which also retains old bonds after adding closer spins. A material model with distance-dependent exchange, a lower bound on site separation across designs, or bounded face aspect ratios would impose additional assumptions. Such a modified optimization problem requires separate analysis.

The counterexample does not affect the proved full-face iterative-triangulation relations or Eq.~\eqref{eq:knownbound}. It provides neither the exact critical temperature of the large fan nor a replacement universal optimum. No claim of historical uniqueness of the decorated graph or transformation method is needed for the strict separation. The comparison concerns the precise published 2026 conjecture; the bounded literature screen does not establish exhaustive worldwide priority.

\section{Reproducibility and AI assistance}
The accompanying source package includes the coordinate and face generator, exact trace and homology checkers, finite radical calculation, complete rational certificates, and regeneration instructions. The primary arithmetic requires Python's standard library; the optional symbolic Kac--Ward and polynomial checks use SymPy. The figure is an original coordinate-based vector drawing. No sampled critical temperature or fitted curve is used.

OpenAI Codex assisted substantively with the construction, literature synthesis, derivations, code generation and debugging, arithmetic checks, figure code, and manuscript preparation. The work was directed toward the specified uniform periodic graph model and exact conjecture. Verification used the distinct representations and published calibrations described above, including a separately written original-edge homology calculation. These checks were computational and AI-assisted; human scientific review and final tool-version disclosure remain pending for this local draft.

\appendix
\section{Finite arithmetic certificates}\label{app:arithmetic}
The ordering inequality~\eqref{eq:integercertificate} is checked by the following complete integer calculation:
\begin{quote}\ttfamily\small
w1, w2 = 11, 10\\
for j in range(1023):\\
\hspace*{1em}w1, w2 = (121*w1 + 100*w2,\\
\hspace*{6em}100*w1 + 100*w2)\\
D = 21 * 210**1023\\
assert 3641*w1 + 1100*w2 > 33000*D
\end{quote}
It verifies a strict rational margin greater than one. No decimal sign test enters the calculation.

For each rational \(x\geq0\), fix \(S=10^{60}\) and compute
\begin{equation}
 k=\operatorname{isqrt}
       \!\left(\left\lfloor\frac{\operatorname{num}(x)S^2}
                                      {\operatorname{den}(x)}\right\rfloor\right).
\end{equation}
Then verify
\begin{equation}
 (k/S)^2\leq x<((k+1)/S)^2.
\end{equation}
Starting with these bounds on \(\sqrt3\), propagate endpoints through Eq.~\eqref{eq:radicals}. Its radicand and the entire map are increasing for \(y>1\). Nine steps give an exact rational upper endpoint smaller than 20. This finite enclosure needs no bound on an asymptotic iteration tail.

For the published temperature calibrations only, logarithms are enclosed by
\begin{equation}
\begin{split}
 2\sum_{k=0}^{m-1}\frac{z^{2k+1}}{2k+1}
 <\ln r\leq
 2\sum_{k=0}^{m-1}\frac{z^{2k+1}}{2k+1}\\
 +\frac{2z^{2m+1}}{(2m+1)(1-z^2)},\qquad
 z=\frac{r-1}{r+1}\in(0,1).
\end{split}
\end{equation}
The positive series and displayed remainder bound control the transcendental conversion. Calibration decimals play no role in the counterexample inequalities.

\section{Original-edge homology elimination}
Represent each old triangular loop by homology label \(1,2,3\), where a translation \((d_x,d_y)\) has label \((d_x\bmod2)+2(d_y\bmod2)\). The \(A\) and \(B\) spokes to \(x_j\) have labels \(0,1\), the \(C x_1\) spoke label \(2\), and chain edges label \(0\).

At rational \(t=p/q\), a frontier state stores hub parity, current internal-site parity, and a homology label. Excluding an edge multiplies its integer weight by \(q\); including it multiplies by \(p\), toggles endpoint parities, and adds its homology. A hub loop toggles hub parity twice. Process the three old loops, then the spokes at each \(x_j\). Its outgoing chain edge must be included exactly when its current parity is odd. This enforces even parity before eliminating the site and initializes the next site's parity. At the last site retain only states with both parities even.

Induction on processed edges shows that each state sums every compatible edge subset exactly once, with every eliminated site even. There are at most sixteen states. After \(3N+3\) edges, the four final weights divided by \(q^{3N+3}\) are the original graph's homology sectors. At \(N=1024,p=1,q=21\), their signed difference is negative and agrees exactly with Eq.~\eqref{eq:originalF}. This supplies a direct original-graph arithmetic check of the thermodynamic sign.

\bibliography{manuscript/references}
\end{document}
